\documentclass[11pt]{article}
\usepackage[a4paper,margin=1in]{geometry}
\usepackage{fontspec}
\usepackage[boldfont=false]{xeCJK}
\setCJKmainfont{FandolSong-Regular.otf}
\usepackage{amsmath}
\usepackage{amssymb}
\usepackage{array}
\usepackage{booktabs}
\usepackage{graphicx}
\usepackage{adjustbox}
\usepackage{enumitem}
\setlist[itemize]{topsep=2pt,partopsep=0pt,parsep=0pt,itemsep=2pt,leftmargin=1.4em}
\setlist[enumerate]{topsep=2pt,partopsep=0pt,parsep=0pt,itemsep=2pt,leftmargin=1.6em}
\usepackage{parskip}
\usepackage{fvextra}
\fvset{breaklines=true,breakanywhere=true,fontsize=\small}
\setlength{\parindent}{0pt}

\begin{document}

\begin{center}
  {\LARGE\bfseries Nucleic acids and protein synthesis}\\[0.35em]
  {\large A Level Biology}
\end{center}
\vspace{0.6em}

\section*{Nucleotides — the building blocks}

\textbf{Nucleic acids} 核酸 (DNA and RNA) are polymers of small units called \textbf{nucleotides} 核苷酸. Each nucleotide is made of three parts joined together:

\begin{itemize}
\item a \textbf{phosphate} 磷酸 group,

\item a sugar (a 5-carbon sugar),

\item a nitrogen-containing \textbf{base} 碱基.

\end{itemize}

\textbf{ATP} is a special nucleotide. It has the base \textbf{adenine} 腺嘌呤, the sugar \textbf{ribose} 核糖, and \textbf{three} phosphate groups. Breaking off the last phosphate releases \textbf{energy} 能量 for the cell.

There are five bases, in two groups:

\begin{itemize}
\item \textbf{purines} 嘌呤 have a \textbf{double ring} (two rings): adenine and \textbf{guanine} 鸟嘌呤.

\item \textbf{pyrimidines} 嘧啶 have a \textbf{single ring} (one ring): \textbf{cytosine} 胞嘧啶, \textbf{thymine} 胸腺嘧啶 and \textbf{uracil} 尿嘧啶.

\end{itemize}

\section*{The structure of DNA}

A DNA molecule is two strands twisted together into a \textbf{double helix} 双螺旋.

Each \textbf{strand} 链 has a backbone of alternating sugar (here the sugar is \textbf{deoxyribose} 脱氧核糖) and phosphate. The sugar of one nucleotide is joined to the phosphate of the next by a \textbf{phosphodiester bond} 磷酸二酯键.

The two strands are held together by their bases, which meet in the middle. The pairing is exact — this is \textbf{complementary base pairing} 碱基互补配对:

\begin{itemize}
\item A always pairs with T, held by \textbf{two} \textbf{hydrogen bonds} 氢键.

\item C always pairs with G, held by \textbf{three} hydrogen bonds (so a C–G \textbf{base pair} 碱基对 is harder to separate).

\end{itemize}

The two strands run in opposite directions: one goes 5′ to 3′ while the other goes 3′ to 5′. We say they are \textbf{antiparallel} 反平行.

\section*{DNA replication}

DNA \textbf{replication} 复制 (copying) happens during the S phase of the cell cycle. It is \textbf{semi-conservative} 半保留复制: each new molecule keeps one old strand and one new strand. The steps are:

\begin{enumerate}
\item the double helix unwinds and the hydrogen bonds break, so the two strands separate.

\item each old strand acts as a template. Free nucleotides pair with the exposed bases by complementary base pairing.

\item the \textbf{enzyme} 酶 \textbf{DNA polymerase} 聚合酶 joins the new nucleotides into a strand. It can only add nucleotides in the 5′ to 3′ direction.

\end{enumerate}

Because of that 5′ to 3′ rule, the two new strands are made differently:

\begin{itemize}
\item the \textbf{leading strand} 前导链 is built continuously, following the unwinding.

\item the \textbf{lagging strand} 后随链 is built in short pieces, working away from the unwinding point. The enzyme \textbf{DNA ligase} 连接酶 then joins these pieces together.

\end{itemize}

\section*{RNA}

RNA is also made of nucleotides, but it is a single strand, its sugar is ribose, and it uses uracil in place of thymine. The most important type here is \textbf{messenger RNA (mRNA)}, which carries a copy of a gene's instructions out of the nucleus to be used.

\section*{The genetic code}

A \textbf{gene} 基因 is a sequence of DNA nucleotides that codes for one \textbf{polypeptide} 多肽.

The code is read in \textbf{triplets} 三联体 — groups of three bases. Each triplet either codes for one specific \textbf{amino acid} 氨基酸, or acts as a start or stop signal. The code is \textbf{universal}: nearly all living things use the same triplets for the same amino acids.

There are 64 possible triplets but only 20 common amino acids, so most amino acids are coded by more than one triplet.

\section*{Protein synthesis: transcription and translation}

\subsection*{Transcription (in the nucleus)}

The DNA gene is copied into mRNA. This is \textbf{transcription} 转录.

\begin{itemize}
\item the strand of DNA that is copied is the \textbf{template strand} 模板链; the partner strand is the \textbf{non-transcribed strand} 非转录链.

\item the enzyme \textbf{RNA polymerase} joins RNA nucleotides that pair with the template bases (with uracil pairing to adenine).

\item in eukaryotes the first RNA made (the \textbf{primary transcript} 初级转录本) contains coding parts called \textbf{exons} 外显子 and non-coding parts called \textbf{introns} 内含子. The introns are cut out and the exons joined together to form the finished mRNA.

\end{itemize}

\subsection*{Translation (at the ribosome)}

The mRNA leaves the nucleus and attaches to a \textbf{ribosome} 核糖体. Building the polypeptide from the mRNA code is \textbf{translation} 翻译.

\begin{itemize}
\item the mRNA is read in \textbf{codons} 密码子 (each codon is one triplet of mRNA bases).

\item molecules of \textbf{transfer RNA (tRNA)} bring amino acids to the ribosome. Each tRNA has an \textbf{anticodon} 反密码子 that pairs with a matching codon.

\item as the codons are read in order, the ribosome joins the amino acids with \textbf{peptide bonds} 肽键, building the polypeptide.

\end{itemize}

\section*{Gene mutations}

A gene \textbf{mutation} 突变 is a change in the base sequence of a DNA molecule. It may change the polypeptide made. There are three types:

\begin{itemize}
\item \textbf{substitution} 替换 — one base is swapped for a different base. This changes at most one amino acid, and sometimes none (because most amino acids have more than one triplet).

\item \textbf{deletion} 缺失 — a base is removed.

\item \textbf{insertion} 插入 — an extra base is added.

\end{itemize}

A deletion or insertion shifts how every later triplet is read, so it usually changes many amino acids after that point and has a large effect on the polypeptide.


\end{document}
